\documentclass[aspectratio=169,12pt]{beamer}

\usepackage{ctex}
\usepackage{amsmath}
\usepackage{amssymb}
\usepackage{graphicx}
\usepackage{booktabs}
\usepackage{colortbl}
\usetheme{Madrid}

\title{多元一次方程组求解}
\subtitle{从消元法到克拉默法则}
\author{}
\date{2026年3月}

\begin{document}

\frame{\titlepage}

\section{问题引入}

\begin{frame}
\frametitle{核心问题}
\textbf{二元方程组}会解，三元、四元甚至更多未知数呢？

\[
\begin{cases}
x + y + z = 6 \\
x - y + z = 4 \\
x + y - z = 2
\end{cases}
\]

\bigskip
\pause
\centerline{\Large 有没有一种方法可以统一求解？}
\end{frame}

\section{消元法}

\begin{frame}
\frametitle{消元法示例}
\textbf{二元例子}：
\[
\begin{cases} x + y = 2 \\ x - y = 0 \end{cases}
\]

消元步骤：
\begin{enumerate}
    \item $(x+y) - (x-y) = 2-0$
    \item $2y = 2 \Rightarrow y = 1$
    \item $x + 1 = 2 \Rightarrow x = 1$
\end{enumerate}
\end{frame}

\begin{frame}
\pause
\textbf{三元例子}：
\[
\begin{cases}
x + y + z = 6 \\
x - y + z = 4 \\
x + y - z = 2
\end{cases}
\]

消元步骤：
\begin{enumerate}
    \item (1)-(2)：$2y = 2 \Rightarrow y = 1$
    \item (1)-(3)：$2z = 4 \Rightarrow z = 2$
    \item 代入(1)：$x + 1 + 2 = 6 \Rightarrow x = 3$
\end{enumerate}
\end{frame}

\begin{frame}
\pause
\textbf{关键发现}：
\centerline{消元法可以推广到任意多元方程组！}

\bigskip
\pause
在消元过程中：
\begin{itemize}
    \item 出现 $0 = 0$ → 可能有无穷多解
    \item 出现 $0 = 1$ → 无解
    \item 正好消元到只剩一个未知数 → 唯一解
\end{itemize}
\end{frame}

\section{消元法的矩阵表示}

\begin{frame}
\fontsize{10}{12}\selectfont
\centerline{\textbf{方程组}}
\[
\begin{cases}
a_{11}x_1 + a_{12}x_2 + a_{13}x_3 = b_1 \\
a_{21}x_1 + a_{22}x_2 + a_{23}x_3 = b_2 \\
a_{31}x_1 + a_{32}x_2 + a_{33}x_3 = b_3
\end{cases}
\]

\bigskip
\centerline{\textbf{矩阵形式}}
\[
A = \begin{pmatrix}
a_{11} & a_{12} & a_{13} \\
a_{21} & a_{22} & a_{23} \\
a_{31} & a_{32} & a_{33}
\end{pmatrix}, \quad
x = \begin{pmatrix} x_1 \\ x_2 \\ x_3 \end{pmatrix}, \quad
b = \begin{pmatrix} b_1 \\ b_2 \\ b_3 \end{pmatrix}
\]

\centerline{简记为：$Ax = b$}
\end{frame}

\begin{frame}
\fontsize{10}{12}\selectfont
\centerline{\textbf{消元法对应初等行变换}}

\begin{center}
\begin{tabular}{|c|c|}
\hline
消元法操作 & 矩阵操作 \\
\hline
方程乘以常数 & 行乘以常数 \\
\hline
方程相加减 & 行相加减 \\
\hline
交换两个方程 & 交换两行 \\
\hline
\end{tabular}
\end{center}

\pause
\bigskip
\centerline{这些操作叫做\textbf{初等行变换}}
\end{frame}

\section{初等行变换的性质}

\begin{frame}
\fontsize{11}{13}\selectfont
\centerline{\textbf{核心性质：初等行变换不改变方程组的解}}

\bigskip
\pause
\centerline{直观理解}
\begin{itemize}
    \item 交换两个方程顺序 → 解不变
    \item 方程乘以非零常数 → 解不变
    \item 方程相加 → 解不变（等量替换）
\end{itemize}
\end{frame}

\begin{frame}
\fontsize{10}{12}\selectfont
\centerline{\textbf{初等行变换与行列式}}

\[
\det\begin{pmatrix} a & b \\ c & d \end{pmatrix} = ad - bc
\]

\begin{center}
\begin{tabular}{|c|c|}
\hline
初等行变换 & 行列式的变化 \\
\hline
交换两行 & 变号（乘以 -1） \\
\hline
行乘以常数 $k$ & 行列式乘以 $k$ \\
\hline
行乘以 $k$ 加到另一行 & \textbf{不变} \\
\hline
\end{tabular}
\end{center}

\pause
\bigskip
\centerline{\textbf{重要}：倍加变换不改变行列式！}
\end{frame}

\section{对角矩阵情况}

\begin{frame}
\fontsize{11}{13}\selectfont
\centerline{\textbf{什么是对角矩阵？}}

\[
D = \begin{pmatrix}
d_1 & & \\
& d_2 & \\
& & \ddots \\
& & & d_n
\end{pmatrix}
\]

\pause
\bigskip
\centerline{\textbf{行列式}：$\det(D) = d_1 \cdot d_2 \cdots d_n$}
\end{frame}

\begin{frame}
\fontsize{11}{13}\selectfont
\centerline{\textbf{对角方程组的解}}

\[
\begin{cases}
d_1 x_1 = b_1 \\
d_2 x_2 = b_2 \\
\vdots \\
d_n x_n = b_n
\end{cases}
\]

\pause
\bigskip
\centerline{\textbf{解}：$x_i = \dfrac{b_i}{d_i}$}

\bigskip
\pause
如果 $d_i = 0$：
\begin{itemize}
    \item 对应 $b_i = 0$ → 无穷多解
    \item 对应 $b_i \neq 0$ → 无解
\end{itemize}
\end{frame}

\section{秩的概念}

\begin{frame}
\fontsize{10}{12}\selectfont
\centerline{\textbf{行列式为零意味着什么？}}

\[
\det(A) = 0 \iff \text{矩阵的"有效信息"不足}
\]

\bigskip
\pause
\centerline{\textbf{从方程组角度看：秩}}

\bigskip
例子：
\[
\begin{cases}
x + y = 2 \\
2x + 2y = 4
\end{cases}
\]

第二个方程是第一个的2倍，\textbf{没有提供新信息}！

\pause
\bigskip
\centerline{\textbf{定义}：不能由其他方程线性组合得到的方程个数}
\end{frame}

\begin{frame}
\fontsize{10}{12}\selectfont
\centerline{\textbf{秩与解的情况}}

\begin{center}
\begin{tabular}{|c|c|c|}
\hline
情况 & 秩 & 解的情况 \\
\hline
秩 = 未知数个数 & $\text{rank}(A) = n$ & 唯一解 \\
\hline
秩 < 未知数个数 & $\text{rank}(A) < n$ & 无穷多解 \\
\hline
增广矩阵秩 > 系数矩阵秩 & $\text{rank}[A|b] > \text{rank}(A)$ & 无解 \\
\hline
\end{tabular}
\end{center}
\end{frame}

\section{克拉默法则}

\begin{frame}
\fontsize{10}{12}\selectfont
\centerline{\textbf{克拉默法则（Cramer's Rule）}}

\bigskip
对于 $n$ 元线性方程组 $Ax = b$，如果 $\det(A) \neq 0$，则：
\[
x_i = \frac{\det(A_i)}{\det(A)}
\]

其中 $A_i$ 是把 $A$ 的第 $i$ 列换成 $b$ 得到的矩阵。
\end{frame}

\begin{frame}
\fontsize{9}{11}\selectfont
\centerline{\textbf{证明（线性列变换）}}

设 $A = (a_1, a_2, ..., a_n)$，方程 $Ax = b$ 写成：
\[
x_1 a_1 + x_2 a_2 + \cdots + x_n a_n = b
\]

\pause
\bigskip
考虑 $A_i = (a_1, ..., a_{i-1}, b, a_{i+1}, ..., a_n)$

\pause
\bigskip
根据行列式对列的线性性质：
\[
\det(A_i) = x_i \cdot \det(A)
\]

\pause
\bigskip
所以：
\[
x_i = \frac{\det(A_i)}{\det(A)}
\]

\centerline{\hfill $\square$}
\end{frame}

\section{结论}

\begin{frame}
\centerline{\textbf{结论}}

\bigskip
\begin{enumerate}
    \item \textbf{通用求解方法}：消元法（初等行变换）
    \item \textbf{矩阵表示}：$Ax = b$
    \item \textbf{行列式性质}：倍加变换不改变行列式
    \item \textbf{解的判断}：$\text{rank}(A)$ vs $\text{rank}[A|b]$
    \item \textbf{公式解}：克拉默法则
\end{enumerate}
\end{frame}

\end{document}
